\section{Related work}
\label{sec:related}

\paragraph{Defragmentation for modern storage.}
FragPicker~\cite{park2021fragpicker} and its journal version~\cite{park2023tocs}
argue that request splitting, not seek distance, is what fragmentation costs on
flash and Optane, and migrate only the data the application's requests straddle.
The artifact's evaluation reproduced for its synthetic workloads; its YCSB and
FileBench experiments are not part of the artifact~\cite{sosp21ae}. ScaleDefrag~\cite{lee2026scaledefrag}
parallelises e4defrag's fragmentation check and data relocation with one thread per
file and reports up to 3.8$\times$ shorter defragmentation on real SATA and Optane
SSDs; there the work per file is data movement, which does benefit from
concurrency. Neither measures the cost of FragPicker's analysis phase, which is our
subject. File-system aging~\cite{smith1997aging,conway2017aging,kadekodi2018geriatrix}
produces the fragmented layouts such tools work on; we only need layouts with
controlled extent counts and use a simpler punch-and-rewrite procedure.

\paragraph{Benchmarking methodology.}
Our comparison follows McSherry et al.'s COST argument~\cite{mcsherry2015cost}: a
scalable implementation should be measured against a competent single-threaded
one. Hoefler and Belli~\cite{hoefler2015scientific} catalogue reporting practice for
parallel speedups, including baselines, repetitions and confidence intervals. For
storage specifically, Traeger et al.~\cite{traeger2008nine} and Tarasov et
al.~\cite{tarasov2011benchmarking} show how strongly results depend on cache state
and setup, and Mytkowicz et al.~\cite{mytkowicz2009wrong} how innocuous-looking setup
choices bias measurements. Our contribution in this space is a concrete case where
the latency model itself, not the measurement, produced the speedup.

\paragraph{Emulating slow storage.}
David~\cite{agrawal2011david} emulates large storage systems by modelling their
runtime rather than storing their data. dm-delay~\cite{kerneldmdelay} and nbdkit's
delay filter~\cite{nbdkit} inject latency below the file system. We use them as two
bounds on device concurrency (any number of requests in flight versus one at a
time) and contrast both with a delay injected in user space.
